\begin{abstract}
A previsão probabilística da vida remanescente de elementos de concreto armado sob carbonatação pode apoiar o planejamento de inspeções e de manutenção preventiva. Este trabalho apresenta um emulador estocástico dependente do tempo que aproxima a distribuição condicional da margem entre o cobrimento e a profundidade carbonatada por uma Distribuição Lambda Generalizada. Seus parâmetros são representados por Expansões em Caos Polinomial em idades discretas, e uma camada de redes neurais permite consultar a distribuição entre essas idades. O procedimento é verificado em um problema idealizado de resistência e solicitação e aplicado a um elemento de concreto armado com variabilidade de resistência, umidade e cobrimento, sob a trajetória de CO$_2$ SSP2-4.5. A partir das distribuições emuladas, constroem-se trajetórias por acoplamento quantílico e estimam-se os tempos até a despassivação e até limiares preventivos. A comparação de margens de zero a 12~mm permite avaliar o efeito do critério de intervenção sobre os prazos, considerando a censura no horizonte de 100 anos. \falta{Completar o resumo com as métricas de acurácia, custo e vida remanescente dos exemplos de durabilidade após a conclusão das simulações.}
\end{abstract}

\textbf{Palavras-chave}: vida \'util remanescente; carbonata\c{c}\~ao; emulador
estoc\'astico; distribui\c{c}\~ao lambda generalizada; confiabilidade dependente do
tempo; concreto armado.

\vspace{0.5em}
\noindent\textbf{Keywords}: remaining useful life; carbonation; stochastic emulator;
generalized lambda distribution; time-dependent reliability; reinforced concrete.
